% ============================================
% MEM-18 (ZH)
% 中文翻译 — SPEC-935-R356
% ============================================
\section*{\textbf{MEM-18} --- 逃跑的尝试}
\markboth{MEM-18 --- 逃跑的尝试}{}

\noindent \textit{1943 —— 差点成功的那一年}

\noindent \textbf{1943 —— 差点成功的那一年}

\vspace{0.5em}\noindent\rule{\textwidth}{0.5pt}\vspace{0.5em}

\noindent 我花了两年计划逃跑。

\noindent 1941年到1943年。日复一日，观察作息。星期一迟到的那位护工。储物间那扇关不严的窗户。后院铁丝网上的洞，被草丛遮着。

\noindent 怪物帮我。

\noindent \textit{“你会成功的，”它说，“你会离开这里。会看见世界。会重新感受阳光。”}

\noindent \textit{“你为什么想让我逃跑？”}

\noindent \textit{“因为在这里面你是渺小的。在外面你可以是巨大的。外面有更多的人。更多的食物。更多的……空间。”}

\noindent 我本该起疑的。

\noindent 怪物从来不为我好。
它总是想要对它\textit{自己}好的东西。

\noindent 但那时我还年轻。1943年，我三十六岁。在收容院里度过了二十六年。我想在没有铁栅栏的中间感受风拂过脸庞。

\vspace{0.5em}\noindent\rule{\textwidth}{0.5pt}\vspace{0.5em}

\noindent 逃跑那一夜，天空干净。星星前所未有地明亮。我等护工做完凌晨两点的巡逻，数到一千，然后起身。

\noindent 牢房的门没有锁——我学会了用藏在脚底的一片金属把它撬开。穿过走廊。寂静。储物间。窗户。开着。

\noindent 我跳了出去。

\noindent 草很高。泥土冰凉。我跑。赤脚踩在地上。铁丝网上的洞。我钻了过去。

\noindent 我在\textit{外面}了。

\noindent 世界很大。黑暗。寂静。但它是开阔的。我望着天空，看见星星，前面没有任何铁栅栏。

\noindent 我哭了。

\noindent 不是悲伤。
是喜悦。

\noindent 二十六年之后，我自由了。

\noindent 十五分钟。

\vspace{0.5em}\noindent\rule{\textwidth}{0.5pt}\vspace{0.5em}

\noindent 我奔跑在土路上时，怪物开口了。

\noindent \textit{“你去哪里？”}

\noindent \textit{“远走。”}

\noindent \textit{“远走去哪里？”}

\noindent \textit{“巴伊亚。听说那里有活干。”}

\noindent \textit{“你打算做什么？”}

\noindent \textit{“干活。吃饭。活着。”}

\noindent \textit{“活着？”它笑了，“你不会活着。你只会求生。”}

\noindent \textit{“我学。”}

\noindent \textit{“你学不会。你是疯人院里养大的野兽。你不会和人交谈。你不会看人的眼睛而不吓到他们。你不知道什么是正常。”}

\noindent \textit{“闭嘴。”}

\noindent \textit{“还有你的眼睛？你觉得会有人雇一个眼睛发黄、在黑暗里发光的男人吗？”}

\noindent 我停下了奔跑。

\noindent 站在路中央，喘着气。

\noindent \textit{“闭嘴。”}

\noindent \textit{“你没有地方可去。外面的世界不想要收容院的前患者。世界想要正常人。你不是正常人。你是患者1,260。这是写在你身上的。”}

\noindent \textit{“我把这个数字从皮肤上撕下来。”}

\noindent \textit{“数字不在皮肤上，若阿金。它在灵魂里。你可以把整张皮撕掉，1,260还在那里。”}

\noindent 我坐在路中央。

\noindent 尘土在我周围扬起。

\noindent 星星继续亮着，美丽，遥远。

\vspace{0.5em}\noindent\rule{\textwidth}{0.5pt}\vspace{0.5em}

\noindent 六个小时后，护工们在同一个地方找到了坐在那里的我。

\noindent 我没有跑。
没有躲。
我只是待在那里，等着。

\noindent \textit{“怎么了？”护工问，“放弃了？”}

\noindent \textit{“没有。”我站起来，“只是发现没有地方可去。”}

\noindent 他们把我带了回去。

\noindent 牢房一模一样。混凝土床。桶。铁栅栏。一切如旧。

\noindent \textit{“欢迎回来，”怪物说。}

\noindent \textit{“闭嘴。”}

\noindent 它笑了。

\noindent \textit{“下次，若阿金，你不要逃跑。你留下。你等待。你书写。总有一天，日记会出去。然后……”}

\noindent \textit{“然后什么？”}

\noindent \textit{“然后逃跑就发生了。但不是你出去。是日记。逃跑属于它。”}

\noindent \textit{“那我呢？”}

\noindent \textit{“你留下。你总是留下。日记才是去看世界的那个。”}\NE{逃跑这一段是若阿金整个日记中最连贯、最线性的部分——传统的叙事结构，有头有尾，与其他碎片的破碎形成对比。这种连贯性暗示了事后的重写，或者说是实体为将要发现它的读者整理叙事的手笔。}

\noindent 它说对了。

\noindent 我再也没有试图逃跑。
但日记逃了。

\noindent 日记替我逃了出去。

\vspace{0.5em}
\noindent\textit{\textbf{↪ 路线:} \rota{MEM-09} \rotasep \rota{MEM-19} \rotasep \rota{MEM-21}}
